\documentclass[10pt,a4paper]{amsart}
\usepackage[margin=19mm]{geometry}
\usepackage{amsmath,amssymb,mathtools}
\usepackage[scr=rsfs,cal=euler]{mathalfa}
\usepackage{xfrac}
\usepackage[unicode,hidelinks,bookmarks=false]{hyperref}
\newcommand{\ie}{i.e.\ }
\newcommand{\cf}{cf.\ }
\newcommand{\resp}{respectively\ }
\newcommand{\Subsection}{\subsection}
\providecommand{\nobreakdash}{\nobreak\hbox{-}}
\setlength{\emergencystretch}{2em}
\multlinegap0pt
\usepackage{bm}
\usepackage{bbm}
\usepackage{dsfont}
\usepackage{xspace}
\usepackage{cancel}
\usepackage{mathtools}
\usepackage{amsthm}
\makeatletter
\@ifundefined{theorem}{\newtheorem{theorem}{Theorem}}{}
\@ifundefined{assume}{\newtheorem{assume}{Assumption}}{}
\@ifundefined{definition}{\newtheorem{definition}[theorem]{Definition}}{}
\@ifundefined{lemma}{\newtheorem{lemma}[theorem]{Lemma}}{}
\@ifundefined{proposition}{\newtheorem{proposition}[theorem]{Proposition}}{}
\@ifundefined{remark}{\newtheorem{remark}[theorem]{Remark}}{}
\makeatother
\newcommand{\RNn}{\mathbb{R}^{N\times n}}
\newcommand{\Rmn}{\mathbb{R}^{m\times n}}
\newcommand{\opnorm}{2\rightarrow2}
\newcommand{\rn}{\mathbb{R}^n}
\newcommand{\st}{\mathcal{S}}\newcommand{\sgn}{\mathrm{sign}}
\title[Generalization analysis of an unfolding network for analysis]{Generalization analysis of an unfolding network for analysis-based Compressed Sensing}
\author{Vicky Kouni, Yannis Panagakis}
\date{Source: arXiv:2303.05582v3 (CC BY 4.0)}
\begin{document}
\maketitle

\noindent\emph{Source and licence.} Generalization analysis of an unfolding network for analysis-based Compressed Sensing. Version arXiv:2303.05582v3 (CC BY 4.0), distributed under the Creative Commons Attribution 4.0 International licence, as recorded in the arXiv licence metadata for this exact version and shown on the abstract page https://arxiv.org/abs/2303.05582v3. Full rights notice: licenses/arxiv-2303.05582-CC-BY-4.0.txt.\par\smallskip

\noindent\textit{Editorial scope.} Counted unit: the Proposition whose source label is \texttt{boundedoutput} in the article (source0.tex lines 426--432, source label \texttt{boundedoutput}), delivered together with the article's own complete original proof of it (source0.tex lines 434--454, one \texttt{proof} environment of 21 lines). No mathematics is written, repaired or re-derived: statement and proof are the article's own text line for line. Required context retained uncounted: layer1 (lines 269--274); assume1 (lines 311--313); assume2 (lines 314--316); frame (lines 326--328); phibound (lines 330--332); invab (lines 410--415). Every definition, display and label of those lines is retained unchanged. Omitted from this package and not redistributed: 1--268, 275--310, 317--325, 329--329, 333--409, 416--425, 433--433, 455--821. Nothing inside a retained excerpt is omitted or altered: every retained body is delivered exactly as the article's lines are written, so no omission receipt of any kind accompanies this package. Citations: the retained excerpts carry no citation command at all, so no bibliography entry of the article is transcribed and no reference is redistributed. Reference adaptation: none was needed and none was made; every reference inside the delivered unit resolves inside it, so no unresolved reference or citation is produced. Printed slips kept literal: no printed word, symbol, hypothesis or number of the counted statement or of its proof was altered. Layout and class: the article's own class is \texttt{elsarticle} with packages including inputenc, babel, amsmath, amssymb, amsthm, booktabs, diagbox, float, balance, textcomp, stfloats, subcaption, framed, algorithm2e; the wrapper is the lane's pinned \texttt{amsart} editorial wrapper at 10pt on A4, which restates the theorem-like environments the retained text needs and adds the article's own macro definitions that the retained text needs (\texttt{\textbackslash RNn}, \texttt{\textbackslash Rmn}, \texttt{\textbackslash opnorm}, \texttt{\textbackslash rn}, \texttt{\textbackslash st}), with extra wrapper packages (bm, bbm, dsfont, xspace, cancel, mathtools). The delivered numbering is the unit's own. Assets: no retained span contains a figure environment, a graphics-inclusion command, a PDF-page inclusion, a table or a TikZ picture; no image file is copied into this package, the package declares no asset dependency and no image file is redistributed with it. Licence: version arXiv:2303.05582v3 of this article is distributed under the Creative Commons Attribution 4.0 International licence, as recorded in the arXiv licence metadata for this exact version and shown on the abstract page https://arxiv.org/abs/2303.05582v3. A full rights notice is delivered at licenses/arxiv-2303.05582-CC-BY-4.0.txt. Characters: every retained excerpt is pure ASCII, so the delivered engine, XeLaTeX with its default Latin Modern text font, reports no missing character.
\begin{align}
    \label{layer1}
    f_1(y) & =I'b+I''\st_{\lambda/\rho}(b),\\
    \label{layerk}
    f_k(v) & =\Tilde{\Theta}v+I'b+I''\st_{\lambda/\rho}(\Theta v+b), \quad k = 2, \hdots, L,
\end{align}

\begin{assume}\label{assume1}
    For an analysis operator $\Phi\in\RNn$ with $N>n$, the matrix $S=\Phi^T\Phi$ is invertible.
\end{assume}

\begin{assume}\label{assume2}
    For an analysis operator satisfying Assumption \ref{assume1}, and for appropriately chosen measurement matrix $A\in\Rmn$ and penalty parameter $\rho>0$, it holds $\rho\|S^{-1}\|_{\opnorm}\|A\|_{\opnorm}<1$.
\end{assume}

\begin{definition}\label{frame}
    We define $\mathcal{F}_{\beta }$ to be the class of redundant analysis operators $\Phi\in\RNn$ for which the associated $S$-operator is invertible and has bounded spectral norm by some $0<\beta<\infty$.
\end{definition}

\begin{remark}\label{phibound}
    The invertibility of $S$ in Definition \ref{frame} implies that the rows of $\Phi$ constitute a frame for $\rn$. Hence, $S$ is a frame operator and for some $0<\alpha\leq\beta<\infty$, it holds $\alpha\leq\|S\|_{\opnorm}\leq\beta$ and $\|\Phi\|_{\opnorm}\leq\sqrt{\beta}$.
\end{remark}

\begin{lemma}[Proof in the supplementary material]\label{invab}
    Let $A\in\mathbb{R}^{n\times n}$ be invertible and $B\in\mathbb{R}^{n\times n}$. For a sub-multiplicative matrix norm $\|\cdot\|$ on $\mathbb{R}^{n\times n}$, if it holds $\|A^{-1}\|\|B\|<1$, then $A+B\in\mathbb{R}^{n\times n}$ is invertible. Moreover, we have
    \begin{equation}
        \|(A+B)^{-1}\|\leq\frac{\|A^{-1}\|}{1-\|A^{-1}\|\|B\|}.
    \end{equation}
\end{lemma}

\begin{proposition}\label{boundedoutput}
    Let $k\in\mathbb{N}$. For any $\Phi\in\mathcal{F}_{\beta }$, with $\mathcal{F}_{\beta }$ as in Definition~\ref{frame}, and arbitrary $\lambda,\,\rho>0$ in the definition of $f_\Phi^k$, we have
    \begin{equation}\label{fbound}
        \|f_\Phi^k(Y)\|_F\leq3\|A\|_{\opnorm}\|Y\|_Fq\sqrt{\beta }\sum_{i=0}^{k-1}3^i(1+2q\rho\beta)^i,
    \end{equation}
where $q=\frac{\rho}{\alpha-\rho\|A^TA\|_{\opnorm}}$, and $\alpha$ and $\beta$ are defined as in Remark~\ref{phibound}.
\end{proposition}

\begin{proof}
    We prove \eqref{fbound} via induction. For $k=1$:
    \begin{equation}\label{f1}
        \|f_\Phi^1(Y)\|_F\leq3\|B\|_F\leq3\|A\|_{\opnorm}\|Y\|_F\sqrt{\beta }\|(A^TA+\rho\Phi^T\Phi)^{-1}\|_{\opnorm},
    \end{equation}
which holds by definition of \eqref{layer1}. The invertibility of $S=\Phi^T\Phi$, along with Assumption \ref{assume2}, Remark \ref{phibound} and Lemma \ref{invab}, imply that
\begin{align}
    \|(A^TA+\rho\Phi^T\Phi)^{-1}\|_{\opnorm}&=\|(A^TA+\rho S)^{-1}\|_{\opnorm}\leq\frac{\rho\|S^{-1}\|_{\opnorm}}{1-\rho\|S^{-1}\|_{\opnorm}\|A^TA\|_{\opnorm}}\notag\\ 
    \label{invs}
    &=\frac{\rho}{\alpha-\rho\|A^TA\|_{\opnorm}}:=q,
\end{align}
where in the last step we used the fact that $\beta^{-1}\leq\|S^{-1}\|_{\opnorm}\leq\alpha^{-1}$, for some $0<\alpha\leq\beta<\infty$. Substituting \eqref{invs} into \eqref{f1} yields $\|f_\Phi^1(Y)\|_F\leq3\|A\|_{\opnorm}\|Y\|_Fq\sqrt{\beta}$. Suppose now that \eqref{fbound} holds for some $k\in\mathbb{N}$. Then, for $k+1$:
\begin{align*}
    \|f_\Phi^{k+1}(Y)\|_F\leq&\|\Tilde{\Theta}\|_{\opnorm}\|f_\Phi^k(Y)\|_F+2\|\Theta\|_{\opnorm}\|f_\Phi^k(Y)\|_F+3\|B\|_F\\
    \leq&3\left((1+2\|W\|_{\opnorm})\|f_\Phi^k(Y)\|_F+\|B\|_F\right)\\
    \leq&3\Bigg((1+2q\rho\beta )\left(3\|A\|_{\opnorm}\|Y\|_Fq\sqrt{\beta}\sum_{i=0}^{k-1}3^i(1+2q\rho\beta )^i\right)\\
    &+\|A\|_{\opnorm}\|Y\|_Fq\sqrt{\beta}\Bigg)\\
    =&3\|A\|_{\opnorm}\|Y\|_Fq\sqrt{\beta}\sum_{i=0}^{k}3^i(1+2q\rho\beta )^i.
\end{align*}
The proof follows.
\end{proof}

\end{document}
